% GENERATED FILE. Edit canonical lessons, then run npm run generate:books.
% canonical-source-hash: fnv1a64:2b844e0ec5bdecdb
% canonical-lessons: SA-C353-sasanam, SA-C353-adhitisthati, SA-C353-rastram, SA-C353-sasakah, SA-C353-adesah

\chapter{Who Governs}
\label{ch:sa-who-governs}
\hlchaptermodality{353}

\begin{chapteropening}
\textbf{By the end of this chapter:} \emph{I can say who rules a country and what orders they give.}

Complete the last lesson of chapter 353: I can say who rules a country and what orders they give.
\end{chapteropening}

\section[\texorpdfstring{\sk{शासनम्} (śāsanam)}{शासनम् (śāsanam)}]{\sk{शासनम्} (śāsanam) — government; rule}

\label{lesson:SA-C353-sasanam}



\begin{quote}
Before the new one: say the Sanskrit for shows favour to, then the Sanskrit for forbearance.
\end{quote}

\subsection*{You'll want to know: \sk{शासनम्}}
\textbf{\sk{शासनम्}} — \emph{śāsanam} — \textquotedblleft{}government; rule\textquotedblright{}.

\begin{quote}
\sk{देशस्य} \sk{शासनं} \sk{कः} \sk{करोति}?
\end{quote}

\begin{quote}
\emph{deśasya śāsanaṃ kaḥ karoti?}
\end{quote}

\begin{quote}
Who runs the country's government?
\end{quote}

A neuter noun.

\begin{cousinweb}
\textbf{\sk{शासनम्}} (\emph{śāsanam}) comes from the root √śās, \textquotedblleft{}to rule, to order\textquotedblright{}; MW gives \emph{śāsana} \textquotedblleft{}government, dominion, rule\textquotedblright{} (MW).
\end{cousinweb}

\begin{grammarlens}[title={What you've built}]
One more word for giving reasons and opinions.
\end{grammarlens}

\subsection*{Your turn}
\begin{itemize}
\raggedright
  \item \emph{Say it:} \emph{śāsanam}
  \item \emph{Say it:} \emph{śāsanam}, once more
  \item \emph{Say it:} \emph{deśasya śāsanaṃ kaḥ karoti?} — Who runs the country's government?
  \item \emph{Recall:} say the Sanskrit for shows favour to, then the Sanskrit for forbearance, then say \emph{śāsanam} again
\end{itemize}

\begin{tcolorbox}[breakable,colback=teal!4,colframe=teal!35!black,title={Before you move on}]
What does \sk{शासनम्} mean? (\textquotedblleft{}Government; rule\textquotedblright{}.) Say it once more.
\end{tcolorbox}
\section[\texorpdfstring{\sk{अधितिष्ठति} (adhitiṣṭhati)}{अधितिष्ठति (adhitiṣṭhati)}]{\sk{अधितिष्ठति} (adhitiṣṭhati) — presides over; governs}

\label{lesson:SA-C353-adhitisthati}



\begin{quote}
Before the new one: say the Sanskrit for forbearance, then the Sanskrit for government.
\end{quote}

\subsection*{You'll want to know: \sk{अधितिष्ठति}}
\textbf{\sk{अधितिष्ठति}} — \emph{adhitiṣṭhati} — \textquotedblleft{}presides over; governs\textquotedblright{}.

\begin{quote}
\sk{मन्त्री} \sk{सभाम्} \sk{अधितिष्ठति।}
\end{quote}

\begin{quote}
\emph{mantrī sabhām adhitiṣṭhati.}
\end{quote}

\begin{quote}
The minister presides over the assembly.
\end{quote}

Past (he …): \sk{सः} \sk{सभाम्} \sk{अधिष्ठितवान्}, \emph{saḥ sabhām adhiṣṭhitavān}.

\begin{cousinweb}
\textbf{\sk{अधितिष्ठति}} (\emph{adhitiṣṭhati}) is built on \emph{adhi}, \textquotedblleft{}over\textquotedblright{}, and the root √sthā, \textquotedblleft{}to stand\textquotedblright{}; MW gives \emph{adhi-√ṣṭhā} \textquotedblleft{}to superintend, to govern\textquotedblright{} (MW).
\end{cousinweb}

\begin{grammarlens}[title={What you've built}]
One more word for giving reasons and opinions.
\end{grammarlens}

\subsection*{Your turn}
\begin{itemize}
\raggedright
  \item \emph{Say it:} \emph{adhitiṣṭhati}
  \item \emph{Say it:} \emph{adhitiṣṭhati}, once more
  \item \emph{Say it:} \emph{mantrī sabhām adhitiṣṭhati.} — The minister presides over the assembly.
  \item \emph{Recall:} say the Sanskrit for forbearance, then the Sanskrit for government, then say \emph{adhitiṣṭhati} again
\end{itemize}

\begin{tcolorbox}[breakable,colback=teal!4,colframe=teal!35!black,title={Before you move on}]
What does \sk{अधितिष्ठति} mean? (\textquotedblleft{}Presides over; governs\textquotedblright{}.) Say it once more.
\end{tcolorbox}
\section[\texorpdfstring{\sk{राष्ट्रम्} (rāṣṭram)}{राष्ट्रम् (rāṣṭram)}]{\sk{राष्ट्रम्} (rāṣṭram) — a nation; a realm}

\label{lesson:SA-C353-rastram}



\begin{quote}
Before the new one: say the Sanskrit for government, then the Sanskrit for presides over.
\end{quote}

\subsection*{You'll want to know: \sk{राष्ट्रम्}}
\textbf{\sk{राष्ट्रम्}} — \emph{rāṣṭram} — \textquotedblleft{}a nation; a realm\textquotedblright{}.

\begin{quote}
\sk{भारतं} \sk{महत्} \sk{राष्ट्रम्} \sk{अस्ति।}
\end{quote}

\begin{quote}
\emph{bhārataṃ mahat rāṣṭram asti.}
\end{quote}

\begin{quote}
India is a great nation.
\end{quote}

A neuter noun.

\begin{cousinweb}
\textbf{\sk{राष्ट्रम्}} (\emph{rāṣṭram}) comes from the root √rāj, \textquotedblleft{}to reign\textquotedblright{}; MW gives \emph{rāṣṭra} \textquotedblleft{}a kingdom, a realm\textquotedblright{} and \textquotedblleft{}a people, a nation\textquotedblright{} (MW).
\end{cousinweb}

\begin{grammarlens}[title={What you've built}]
One more word for giving reasons and opinions.
\end{grammarlens}

\subsection*{Your turn}
\begin{itemize}
\raggedright
  \item \emph{Say it:} \emph{rāṣṭram}
  \item \emph{Say it:} \emph{rāṣṭram}, once more
  \item \emph{Say it:} \emph{bhārataṃ mahat rāṣṭram asti.} — India is a great nation.
  \item \emph{Recall:} say the Sanskrit for government, then the Sanskrit for presides over, then say \emph{rāṣṭram} again
\end{itemize}

\begin{tcolorbox}[breakable,colback=teal!4,colframe=teal!35!black,title={Before you move on}]
What does \sk{राष्ट्रम्} mean? (\textquotedblleft{}A nation; a realm\textquotedblright{}.) Say it once more.
\end{tcolorbox}
\section[\texorpdfstring{\sk{शासकः} (śāsakaḥ)}{शासकः (śāsakaḥ)}]{\sk{शासकः} (śāsakaḥ) — a ruler; a governor}

\label{lesson:SA-C353-sasakah}



\begin{quote}
Before the new one: say the Sanskrit for presides over, then the Sanskrit for a nation.
\end{quote}

\subsection*{You'll want to know: \sk{शासकः}}
\textbf{\sk{शासकः}} — \emph{śāsakaḥ} — \textquotedblleft{}a ruler; a governor\textquotedblright{}.

\begin{quote}
\sk{शासकः} \sk{देशस्य} \sk{सुखं} \sk{चिन्तयेत्।}
\end{quote}

\begin{quote}
\emph{śāsakaḥ deśasya sukhaṃ cintayet.}
\end{quote}

\begin{quote}
A ruler should think of the country's well-being.
\end{quote}

A masculine noun.

\begin{cousinweb}
\textbf{\sk{शासकः}} (\emph{śāsakaḥ}) comes from the root √śās, \textquotedblleft{}to rule\textquotedblright{}, with the agent ending \emph{-aka}; MW gives \emph{śāsaka} \textquotedblleft{}a chastiser, a teacher, a governor, a ruler\textquotedblright{} (MW).
\end{cousinweb}

\begin{grammarlens}[title={What you've built}]
One more word for giving reasons and opinions.
\end{grammarlens}

\subsection*{Your turn}
\begin{itemize}
\raggedright
  \item \emph{Say it:} \emph{śāsakaḥ}
  \item \emph{Say it:} \emph{śāsakaḥ}, once more
  \item \emph{Say it:} \emph{śāsakaḥ deśasya sukhaṃ cintayet.} — A ruler should think of the country's well-being.
  \item \emph{Recall:} say the Sanskrit for presides over, then the Sanskrit for a nation, then say \emph{śāsakaḥ} again
\end{itemize}

\begin{tcolorbox}[breakable,colback=teal!4,colframe=teal!35!black,title={Before you move on}]
What does \sk{शासकः} mean? (\textquotedblleft{}A ruler; a governor\textquotedblright{}.) Say it once more.
\end{tcolorbox}
\section[\texorpdfstring{\sk{आदेशः} (ādeśaḥ)}{आदेशः (ādeśaḥ)}]{\sk{आदेशः} (ādeśaḥ) — an official order; a command}

\label{lesson:SA-C353-adesah}



\begin{quote}
Before the new one: say the Sanskrit for a nation, then the Sanskrit for a governor.
\end{quote}

\subsection*{You'll want to know: \sk{आदेशः}}
\textbf{\sk{आदेशः}} — \emph{ādeśaḥ} — \textquotedblleft{}an official order; a command\textquotedblright{}.

\begin{quote}
\sk{मन्त्री} \sk{आदेशं} \sk{पालयति।}
\end{quote}

\begin{quote}
\emph{mantrī ādeśaṃ pālayati.}
\end{quote}

\begin{quote}
The minister carries out the order.
\end{quote}

A masculine noun.

\begin{cousinweb}
\textbf{\sk{आदेशः}} (\emph{ādeśaḥ}) is built on \emph{ā} and the root √diś, \textquotedblleft{}to point out\textquotedblright{}; MW gives \emph{ā-deśa} \textquotedblleft{}rule, command, order\textquotedblright{} (MW).
\end{cousinweb}

\begin{grammarlens}[title={What you've built}]
One more word for giving reasons and opinions.
\end{grammarlens}

\subsection*{Your turn}
\begin{itemize}
\raggedright
  \item \emph{Say it:} \emph{ādeśaḥ}
  \item \emph{Say it:} \emph{ādeśaḥ}, once more
  \item \emph{Say it:} \emph{mantrī ādeśaṃ pālayati.} — The minister carries out the order.
  \item \emph{Recall:} say the Sanskrit for a nation, then the Sanskrit for a governor, then say \emph{ādeśaḥ} again
\end{itemize}

\begin{tcolorbox}[breakable,colback=teal!4,colframe=teal!35!black,title={Before you move on}]
What does \sk{आदेशः} mean? (\textquotedblleft{}An official order; a command\textquotedblright{}.) Say it once more.
\end{tcolorbox}
